\chapter{Einleitung und Zielstellung}

Die Einleitung ist das erste inhaltliche Kapitel und hat eine klare Aufgabe: Sie führt den Leser an das Thema heran, beschreibt das Problem und formuliert die Ziele der Arbeit. Ein guter Einstieg beginnt allgemein und wird schrittweise konkreter — vom breiten Themenfeld bis zur spezifischen Forschungsfrage.

Wichtig: Die Einleitung enthält in der Regel noch keine eigenen Ergebnisse. Sie motiviert die Arbeit und gibt dem Leser eine Orientierung, was ihn erwartet.

\section{Motivation und Problemstellung}

In diesem Abschnitt wird dargelegt, warum das Thema relevant ist. Dabei hilft es, auf aktuelle Entwicklungen in Industrie oder Forschung zu verweisen. Die Motivation sollte so formuliert sein, dass auch fachfremde Leser den Kern des Problems verstehen.

Abkürzungen wie \ac{KI} oder \ac{FAPS} werden beim ersten Auftreten im Text automatisch ausgeschrieben. Bei jeder weiteren Verwendung erscheint nur noch das Kürzel: \ac{KI}.

\section{Zielsetzung der Arbeit}

Die Zielsetzung beantwortet die Frage: Was soll am Ende der Bearbeitungszeit konkret erreicht sein? Ziele sollten möglichst spezifisch und überprüfbar formuliert werden — nicht \glqq das Thema untersuchen\grqq, sondern z.\,B. \glqq einen Prototyp entwickeln und anhand von drei Kriterien bewerten\grqq.

Optional kann hier auch der Aufbau der Arbeit kurz skizziert werden, damit der Leser weiß, welches Kapitel welchen Beitrag leistet.